\documentclass[11pt]{article}
\usepackage[a4paper,margin=1in]{geometry}
\usepackage{amsmath,amssymb,amsthm,mathtools}
\usepackage{enumitem}
\usepackage{microtype}
\usepackage[hidelinks]{hyperref}

\newtheorem{theorem}{Theorem}[section]
\newtheorem{lemma}[theorem]{Lemma}
\newtheorem{proposition}[theorem]{Proposition}
\newtheorem{external}[theorem]{Published input}
\theoremstyle{definition}
\newtheorem{problem}[theorem]{Problem}
\theoremstyle{remark}
\newtheorem{remark}[theorem]{Remark}

\newcommand{\R}{\mathbb{R}}
\newcommand{\lap}{\Delta}
\newcommand{\1}{\mathbf{1}}

\title{A Counterexample to Nikliborc Problem 128\\with Explicit Topological Control}
\author{}
\date{Revised standalone draft, September 2026}

\newcommand{\publicationstatus}{%
\begin{center}
\fbox{\begin{minipage}{0.92\linewidth}
\small
\textbf{AI EXPLORATORY MATHEMATICAL NOTE}\\[2pt]
\textbf{Status:} E1 - AI cross-checked; \textbf{NO INDEPENDENT HUMAN REVIEW}\\
\textbf{Version:} v1.0\\
\textbf{First public release:} PENDING\\
\textbf{Related Lean formalization:} the article-specific deduction is kernel-checked relative to explicit \texttt{PublishedTheory} assumptions, which are not themselves formalized here; see the end matter.\\
\end{minipage}}
\end{center}
\medskip
}

\hypersetup{pdfauthor={}}

\begin{document}
\maketitle

\publicationstatus

\begin{abstract}
Nikliborc Problem 128 asks whether a homogeneous three-dimensional solid whose Newtonian potential is polynomial throughout the solid and its boundary must be an ellipsoid, under the topological hypothesis appearing in the printed problem. We construct a counterexample satisfying that hypothesis.
\end{abstract}

\begin{quote}
\small
\textbf{Feedback.} This note has not undergone independent human mathematical review. Readers who know relevant prior literature, identify an error, or wish to check the argument are especially encouraged to contact the coordinators listed at the end of the paper or to open an issue in the repository.
\end{quote}


\section{The problem and normalization}

For a bounded measurable body $\Omega\subset\R^3$ define
\begin{equation}\label{eq:newton}
 U_\Omega(x)=\frac1{4\pi}\int_\Omega\frac{dy}{|x-y|}.
\end{equation}
The printed Problem 128 permits a polynomial of arbitrary degree and asks, in particular, whether the stated topological hypothesis together with polynomiality of the Newtonian potential on the solid and its boundary forces the solid to be an ellipsoid \cite{ScottishBook}.

\begin{problem}[Nikliborc 128, modernized]\label{prob:128}
Must every homogeneous unicoherent solid in $\R^3$ whose Newtonian potential is polynomial throughout the solid and its boundary be an ellipsoid?
\end{problem}

The previous draft invoked the fact that Yuan--Liu's small-parameter examples are ``quasi-ellipsoidal.'' That phrase by itself is not a topological proof. We therefore use a near-ball realization whose closed body is explicitly the image of a closed ball under a diffeomorphism.

\section{Published inputs}

\begin{external}[Ball potential]\label{ext:ball}
For $B_a=\{x:|x|<a\}$,
\[
 U_{B_a}(x)=
 \begin{cases}
 \displaystyle\frac{a^2}{2}-\frac{|x|^2}{6},&|x|\le a,\\[5pt]
 \displaystyle\frac{a^3}{3|x|},&|x|\ge a.
 \end{cases}
\]
In particular, with $q_a(x)=a^2/2-|x|^2/6$,
\begin{equation}\label{eq:ballgap}
 U_{B_a}=q_a\ \hbox{on }\overline{B_a},
 \qquad
 U_{B_a}(x)-q_a(x)=\frac{(|x|-a)^2(|x|+2a)}{6|x|}>0
 \ \hbox{for }|x|>a.
\end{equation}
\end{external}

\begin{external}[Obstacle problem and free-boundary stability]\label{ext:obstacleSS}
We use standard whole-space obstacle-problem existence and $C^{1,1}$ regularity together with Theorem 1.1 of Serfaty--Serra \cite{SerfatySerra}. In dimension $3$, after verifying their asymptotic and compact-support hypotheses and the quantitative conditions
\[
 \lap h_0\le-\rho\quad\hbox{near the initial coincidence set},
 \qquad
 u_0-h_0\ge\rho\quad\hbox{away from it},
\]
as well as a uniform two-sided tangent-ball condition on the initial free boundary, the theorem supplies a family of diffeomorphisms close to the identity carrying the initial noncoincidence set and free boundary to the perturbed ones.
\end{external}

\begin{external}[Newtonian representation and continuity]\label{ext:newtonunique}
If $u(x)\to0$ at infinity and $-\lap u=\1_K$ distributionally for a bounded body $K$, then $u=U_K$. We also use the standard continuity of $U_K$ on $\R^3$ for bounded volume density.
\end{external}

\begin{external}[Interior potential of an ellipsoid]\label{ext:ellipsoid}
The Newtonian potential of a homogeneous ellipsoid agrees on its interior with a polynomial of degree at most two.
\end{external}

The ball and ellipsoid facts are classical; see MacMillan \cite{MacMillan}, Karp \cite{Karp}, or Di Fratta \cite{DiFratta}. We also use the standard topological facts that a closed ball is unicoherent and that unicoherence is preserved by homeomorphism.

\section{A topologically controlled polynomial inclusion}

For concreteness choose the cubic polynomial
\begin{equation}\label{eq:H}
 H(x_1,x_2,x_3)=x_1^3-3x_1x_2^2.
\end{equation}
It is nonzero, homogeneous of degree $3$, and harmonic because
\[
 \frac{\partial^2 H}{\partial x_1^2}=6x_1,
 \qquad
 \frac{\partial^2 H}{\partial x_2^2}=-6x_1,
 \qquad
 \frac{\partial^2 H}{\partial x_3^2}=0.
\]
Set
\begin{equation}\label{eq:pt}
 p_t(x)=\frac12-\frac{|x|^2}{6}+tH(x).
\end{equation}
Then $-\lap p_t=1$.

\begin{lemma}[Near-ball polynomial inclusion]\label{lem:nearball}
There exists $\tau>0$ such that for every $|t|<\tau$ there is a bounded domain $K_t\subset\R^3$ with $C^1$ boundary and a diffeomorphism $\Psi_t$ close to the identity for which
\begin{equation}\label{eq:topology}
 \overline{K_t}=\Psi_t(\overline{B_1}),
 \qquad
 \partial K_t=\Psi_t(S^2),
\end{equation}
and
\begin{equation}\label{eq:potentialinside}
 U_{K_t}(x)=p_t(x)\qquad(x\in K_t).
\end{equation}
\end{lemma}

\begin{proof}
Choose
\[
 1<b<r<R
\]
and $\chi\in C_c^\infty(B_R)$ with $0\le\chi\le1$ and $\chi\equiv1$ on $B_r$. Fix $M>0$ and define
\begin{equation}\label{eq:obstacle}
 h_t(x)=\chi(x)p_t(x)-(1-\chi(x))M.
\end{equation}
Let $u_t$ be the whole-space obstacle solution with asymptotic constant $c(t)\equiv0$.

At $t=0$ take $u_0=U_{B_1}$. By \eqref{eq:ballgap}, and because $-M<U_{B_1}$, the closed coincidence set is exactly
\begin{equation}\label{eq:contact}
 C_0:=\{u_0=h_0\}=\overline{B_1};
\end{equation}
the strict inequality holds at every point outside the closed unit ball.

We verify the quantitative stability assumptions rather than hiding them in the phrase ``quasi-ellipsoidal.'' First,
\[
 h_t-h_0=t\chi H,
\]
so the perturbation and its Laplacian are compactly supported and tend to zero at infinity. Moreover $h_t\to-M<0=c(t)$ at infinity. Take $U=B_b$. Since $b<r$, one has $h_0=\tfrac12-|x|^2/6$ on a neighborhood of $\overline U$, hence $\lap h_0=-1$ there. The continuous gap $u_0-h_0$ is strictly positive on the compact annulus $\overline{B_R}\setminus B_b$, while outside $B_R$ it equals $U_{B_1}+M\ge M$. Thus it has a positive lower bound on $\R^3\setminus U$. The sphere $S^2$ has a uniform two-sided tangent-ball radius. Choosing one positive $\rho$ below these three bounds verifies the quantitative hypotheses of Published Input~\ref{ext:obstacleSS}.

The published theorem therefore provides a global diffeomorphism $\Psi_t$ carrying the initial noncoincidence set $\R^3\setminus\overline{B_1}$ and its boundary to their perturbed counterparts. Taking complements gives the perturbed closed coincidence set
\[
 C_t=\Psi_t(\overline{B_1}).
\]
Define $K_t=\operatorname{int}C_t=\Psi_t(B_1)$. Then \eqref{eq:topology} follows. Because the diffeomorphism fixes the complement of $U=B_b$, the coincidence set remains in $B_b\Subset B_r$, where $\chi=1$. Hence on $K_t$,
\[
 u_t=h_t=p_t.
\]
Since $-\lap p_t=1$ on $K_t$ and $u_t$ is harmonic on the noncoincidence set, $C^{1,1}$ regularity excludes a singular boundary measure and yields
\[
 -\lap u_t=\1_{K_t}
\]
distributionally. Published Input~\ref{ext:newtonunique} gives $u_t=U_{K_t}$, proving \eqref{eq:potentialinside}.
\end{proof}

\begin{lemma}[Extension to the boundary]\label{lem:boundary}
For every $|t|<\tau$,
\begin{equation}\label{eq:boundarypotential}
 U_{K_t}(x)=p_t(x)\qquad(x\in\overline{K_t}).
\end{equation}
\end{lemma}

\begin{proof}
Both sides of \eqref{eq:potentialinside} are continuous on $\R^3$: $p_t$ is a polynomial and $U_{K_t}$ is continuous by Published Input~\ref{ext:newtonunique}. Since $K_t$ is dense in its closure, the equality extends to $\overline{K_t}$.
\end{proof}

\section{The high-degree obstruction to ellipsoids}

\begin{lemma}[The cubic term cannot cancel]\label{lem:degree}
If $t\ne0$, then $p_t$ has degree exactly $3$.
\end{lemma}

\begin{proof}
The degree-$3$ homogeneous component of $p_t$ is $tH$. It is nonzero because $t\ne0$ and $H\ne0$, while the remaining terms have degree at most two.
\end{proof}

\begin{lemma}[The near-ball body is not an ellipsoid]\label{lem:nonellipsoid}
For every sufficiently small $t\ne0$, $K_t$ is not an ellipsoid.
\end{lemma}

\begin{proof}
Assume $K_t$ is an ellipsoid. By Published Input~\ref{ext:ellipsoid}, there is a polynomial $e$ of degree at most two such that
\[
 U_{K_t}(x)=e(x)\qquad(x\in K_t).
\]
By Lemma~\ref{lem:nearball}, the same potential equals $p_t$ on the nonempty open set $K_t$. Hence $p_t-e$ vanishes on a nonempty open subset of $\R^3$. The polynomial identity principle implies $p_t=e$ as polynomials. This contradicts Lemma~\ref{lem:degree}.
\end{proof}

\begin{theorem}[Counterexample to Problem 128]\label{thm:counterexample128}
The printed form of Nikliborc Problem 128 admits the following counterexample.
\end{theorem}

\begin{proof}
Choose a sufficiently small nonzero $t$ and set
\[
 T=\overline{K_t}.
\]
By \eqref{eq:topology}, $T$ is homeomorphic to the closed unit ball. Hence it satisfies the historical unicoherence hypothesis. Since $\partial K_t$ is a $C^1$ surface, it has volume measure zero, so taking the homogeneous solid to be $K_t$ or its closure gives the same Newtonian volume potential. Lemma~\ref{lem:boundary} shows that this potential equals the polynomial $p_t$ at every point of the solid and its boundary. Lemma~\ref{lem:nonellipsoid} shows that the body is not an ellipsoid. Thus all hypotheses of Problem~\ref{prob:128} are met while its proposed conclusion fails.
\end{proof}

\section{Relation to Yuan--Liu and provenance}

Yuan and Liu introduced polynomial inclusions and proved the existence of $p$-inclusions of degree greater than two by variational-inequality methods; their Theorem 3.2 gives a local small-parameter existence result \cite{YuanLiu}. Thus the underlying existence phenomenon used to negate Problem 128 is not a new discovery of this note. The role of Lemma~\ref{lem:nearball} is narrower: it uses a smooth near-ball obstacle and the published Serfaty--Serra stability theorem to make the topological connection to the printed Scottish Book hypothesis explicit, rather than treating the word ``quasi-ellipsoidal'' as a topological theorem.


\section*{Acknowledgments}
Chuyang Chen and Chenhao Bian are grateful to their advisor, Fei Yu, for his guidance on this project.

\section*{Independent review and Lean verification notice}
\begingroup
\small
\begin{center}
\begin{minipage}{0.94\linewidth}
\centering
\textbf{\large No independent human mathematical review has been performed.}
\end{minipage}
\end{center}

\medskip
\noindent The accompanying Lean file compiles under Lean 4.33.1 / Mathlib version v4.33.1 (exact mathlib revision \texttt{0df444a360eaa60ab8c11dca51a86af692955474}).

\medskip
\noindent The principal Lean theorem is conditional on the explicit \texttt{PublishedTheory} interface; those external mathematical inputs are not proved in the Lean file.

\medskip
\noindent\textbf{Successful Lean verification establishes correctness of the formalized statement relative to its assumptions. It does not by itself establish that the formalized statement is exactly equivalent to the historical problem as originally formulated.}
\endgroup

\section*{Provenance and contacts}
\begingroup\small\sloppy
\begin{description}
\item[Project curator.] Fei Yu.
\item[AI-generation coordinators and contacts.] Chuyang Chen (\href{mailto:22535018@zju.edu.cn}{22535018@zju.edu.cn}); Chenhao Bian (\href{mailto:bch26@mails.tsinghua.edu.cn}{bch26@mails.tsinghua.edu.cn}).
\item[Mathematical draft generation.] Mathematical drafts were generated with GPT-5.6 Sol under their supervision.
\item[Lean code generation.] The accompanying Lean source was generated and revised with GPT-5.6 Sol under their supervision. This is a code-generation provenance statement and is separate from mathematical review.
\item[Lean build/kernel execution.] The local build and kernel-check commands reported below were executed by Chuyang Chen in the recorded project environment.
\item[Human mathematical verification.] NONE.
\item[Statement-correspondence audit.] NONE. No independent human audit has certified that the natural-language statement and the formal Lean statement are exactly equivalent.
\end{description}
\medskip
\noindent Comments, corrections, historical references, and independent verification are very welcome. For long-term correspondence, readers are encouraged to use the repository issue tracker as well as the contacts above.
\endgroup

\section*{Lean formal verification record}

\begingroup\small\sloppy
This record separates the natural-language statement, the formal Lean statement, kernel checking of the proof term, and the separate audit of the correspondence between the first two. Kernel acceptance is not treated as human verification of the original problem.

\begin{description}
\item[Formalization status.] Conditional F1 for the formal theorem; full-paper/original-problem formalization remains PARTIAL.
\item[Lean source.] \texttt{\detokenize{PaperFormalization/Scottish 128/Main.lean}}.
\item[Principal formal theorem.] \texttt{\detokenize{Nikliborc128.negative_answer_128}} with explicit parameter \texttt{\detokenize{pub : PublishedTheory}}.
\item[Build command.] \texttt{\detokenize{lake env lean "PaperFormalization/Scottish 128/Main.lean"}}.
\item[Build/kernel result.] PASS for the conditional theorem stated above.
\item[Lean version.] Lean 4.33.1 (commit 819816b2e0a3bf405af45ae5c7af2491d8f5bee6).
\item[Library snapshot.] mathlib v4.33.1, exact revision 0df444a360eaa60ab8c11dca51a86af692955474.
\item[Project commit.] PENDING -- the final verified working tree has not yet been committed.
\item[Kernel axiom audit.] The principal theorem depends only on \texttt{propext}, \texttt{Classical.choice}, and \texttt{Quot.sound} at the Lean axiom level. There is no user-defined Lean axiom in that dependency audit.
\item[Explicit mathematical assumptions.] The parameter \texttt{PublishedTheory} contains 27 explicit theorem interfaces. They cover: polynomial/analytic Laplacian bridges (\texttt{laplacian\_eval\_cubic}, \texttt{laplacian\_congr\_nhds}); smooth cutoffs; the homogeneous-ball Newtonian potential (inside formula, exterior gap, positivity, $C^{1,1}$ regularity, decay, and Poisson identity); whole-space obstacle complementarity, existence, uniqueness, and the Serfaty--Serra stability interface; local-integrability and distributional-Laplacian bridges; Newtonian representation/uniqueness, continuity, and kernel integrability; two-sided ball geometry, nullity of $C^1$ sphere images, null-set almost-everywhere facts, homeomorphic images of ball closure/frontier/interior, and unicoherence; and the classical theorem that the interior Newtonian potential of a solid ellipsoid is quadratic. These interfaces are assumptions of the conditional theorem and are not proved inside this Lean file.
\item[Scope note.] The article-specific deductions from \texttt{PublishedTheory} are kernel-checked, but the 27 external theorem interfaces themselves are not formalized here. Consequently a clean \texttt{\#print axioms} result does not make this an unconditional full-paper formalization.
\end{description}

\noindent\textbf{Interpretation of the label.} F1 applies to the conditional formal theorem named \texttt{\detokenize{negative_answer_128}}, with explicit parameter \texttt{\detokenize{pub : PublishedTheory}}. The full mathematical argument remains PARTIAL until the \texttt{PublishedTheory} inputs are themselves formalized (or otherwise covered by a separately declared verification policy). F2 would additionally require a human statement-correspondence audit.
\endgroup

\section{Formalization boundary for Lean}

For the Lean development, public analytic theorems may be represented by accurately stated external theorem interfaces. Every construction and deduction specific to this note must be proved. The near-ball lemma may be implemented once in a shared module with the formalization of Problem 129 and imported here, but it may not be postulated.

\subsection*{Published interfaces allowed as external inputs}
\begin{enumerate}[label=\textbf{E128.\arabic*.}]
\item The ball-potential formula in Published Input~\ref{ext:ball}.
\item Standard whole-space obstacle existence/complementarity and $C^{1,1}$ regularity, together with Serfaty--Serra Theorem 1.1 in the precise form used in Published Input~\ref{ext:obstacleSS}.
\item Newtonian representation/uniqueness and continuity, Published Input~\ref{ext:newtonunique}.
\item The classical ellipsoid interior-potential theorem, Published Input~\ref{ext:ellipsoid}.
\item Standard existence of smooth cutoffs and the standard topological facts concerning the closed ball and unicoherence.
\end{enumerate}

\subsection*{Claims that must be proved in Lean}
\begin{enumerate}[label=\textbf{L128.\arabic*.}]
\item Define the explicit cubic $H$ in \eqref{eq:H} and prove it is nonzero, homogeneous of degree $3$, and harmonic.
\item Define $p_t$ and the smooth obstacle $h_t$, and prove the Laplacian and perturbation identities used in Lemma~\ref{lem:nearball}.
\item Prove the exact initial coincidence-set identity $C_0=\overline{B_1}$.
\item Verify the Serfaty--Serra compact-support, asymptotic, negative-Laplacian, positive-gap, fixed-$c(t)$ and two-sided tangent-ball hypotheses for the concrete family.
\item From the published diffeomorphism theorem, construct $K_t$, prove $\overline{K_t}\simeq\overline{B_1}$, derive the distributional equation $-\lap u_t=\1_{K_t}$, and use the published uniqueness theorem to prove $U_{K_t}=p_t$ on $K_t$.
\item Prove Lemma~\ref{lem:boundary} by continuity, so polynomiality holds on the boundary as required by the printed problem.
\item Prove Lemma~\ref{lem:degree} and the polynomial identity principle step used in Lemma~\ref{lem:nonellipsoid}; nonellipsoidality may not be assumed as a certificate field.
\item Prove that the homeomorphism with a closed ball supplies the historical topological hypothesis and assemble the exact counterexample in Theorem~\ref{thm:counterexample128}.
\end{enumerate}

\begin{thebibliography}{9}
\bibitem{ScottishBook}
R. D. Mauldin (ed.), \emph{The Scottish Book: Mathematics from the Scottish Caf\'e, with Selected Problems from the New Scottish Book}, 2nd ed., Birkh\"auser/Springer, 2015. DOI: 10.1007/978-3-319-22897-6.

\bibitem{YuanLiu}
T. Yuan and L. Liu, ``Polynomial inclusions: definitions, applications, and open problems,'' \emph{Journal of the Mechanics and Physics of Solids} 181 (2023), 105440. DOI: 10.1016/j.jmps.2023.105440.

\bibitem{SerfatySerra}
S. Serfaty and J. Serra, ``Quantitative stability of the free boundary in the obstacle problem,'' \emph{Analysis \& PDE} 11 (2018), no. 7, 1803--1839. DOI: 10.2140/apde.2018.11.1803.

\bibitem{MacMillan}
W. D. MacMillan, \emph{The Theory of the Potential}, McGraw--Hill, New York and London, 1930.

\bibitem{Karp}
L. Karp, ``On the Newtonian potential of ellipsoids,'' \emph{Complex Variables, Theory and Application} 25 (1994), 367--371. DOI: 10.1080/17476939408814757.

\bibitem{DiFratta}
G. Di Fratta, ``The Newtonian potential and the demagnetizing factors of the general ellipsoid,'' \emph{Proceedings of the Royal Society A} 472 (2016), 20160197. DOI: 10.1098/rspa.2016.0197.
\end{thebibliography}

\end{document}
